\section{本章小结}
\label{sec:cons-summary}

本章从车流的收支平衡出发，研究传播速度依赖未知量时会发生什么。特征线是直线，斜率由初值决定；只要相邻特征线的速度差使$1+a't$为零，斜率就在有限时间发散。此后解的值仍然有界，但$C^1$解不再存在，初值~\eqref{eq:cons-bump}的破裂时刻是$T^{*}=8\sqrt{3}/9$，斜坡~\eqref{eq:cons-ramp}的破裂时刻是$T^{*}=1$。

收支平衡不依赖导数，于是弱形式~\eqref{eq:cons-weak}接纳了间断，Rankine--Hugoniot条件~\eqref{eq:cons-rh}把间断的速度变成了通量的弦斜率。但这只是总量守恒的要求。同一阶跃初值有一整族弱解~\eqref{eq:cons-weak-family}，零初值也有非零的弱解~\eqref{eq:cons-zero-data}，守恒本身选不出解。

熵条件补上了这一步。从特征线看，它要求特征线汇入间断；从小黏性看，它是黏性耗散在极限中留下的不等式；从通量图像看，它是弦与图像的上下关系。Riemann问题把这些条件落实为计算：严格凸通量下，$u_l>u_r$得到激波，$u_l<u_r$得到稀疏波。表~\ref{tab:cons-summary}把各个解的概念放在一起。

\begin{BookInlineTable}
  \centering
  \caption{标量守恒律的解的概念及其适用条件。}
  \label{tab:cons-summary}
  \begin{tabular}{p{16mm}p{32mm}p{24mm}p{20mm}}
    \toprule
    概念 & 要求 & 存在范围 & 唯一性\\
    \midrule
    $C^1$解 & 逐点满足方程 & 到$T^{*}$为止 & 有\\
    弱解 & 式~\eqref{eq:cons-weak}，间断满足跳跃条件 & 整体（含熵解） & 无\\
    熵解 & 弱解加所有凸熵不等式 & 整体，$u_0\in L^{\infty}$ & 有，且$L^1$收缩\\
    \bottomrule
  \end{tabular}
\end{BookInlineTable}

最后，变量加倍把熵不等式变成两个熵解之间的比较，给出局部$L^1$收缩~\eqref{eq:cons-local-contraction}，由此得到唯一性、比较原理和全变差不增。这些结论依赖通量的正则性和解的有界性，$L^1$距离还要求初值之差可积；存在性依赖消失黏性或差分格式的极限，本章未证。一维标量以外的情形，尤其是守恒律系统，不能照搬这些结论。

守恒和正确的解选择缺一不可：只有守恒，解太多；只有特征线，解在相交之后不存在。第~\ref{chap:pde-numerical-methods}章的有限体积方法在每个单元界面上解一个Riemann问题，复用的正是这里的局部传播结构。
